\documentclass[11pt]{article}
\usepackage[margin=1in]{geometry}
\usepackage{amsmath,amssymb,amsthm,mathtools}
\usepackage{graphicx}
\usepackage{microtype}
\usepackage[hidelinks]{hyperref}

\newtheorem{theorem}{Theorem}
\newtheorem{lemma}{Lemma}
\newcommand{\Rbd}{\mathcal R}
\newcommand{\M}{\mathcal M}
\newcommand{\PhiClass}{\boldsymbol\Phi}
\newcommand{\E}{\mathbb E}

\title{A quantitative-aperture reduction from nontangential RMF to dyadic RMF}
\author{Automated mathematics run \texttt{fa\_banach\_001}}
\date{21 July 2026}

\begin{document}
\maketitle

\noindent\textbf{Status.} Conditional reduction, likely valid, for expert review.  The
argument below proves the stated implication from an explicit quantitative aperture
hypothesis.  It does \emph{not} prove that ordinary nontangential RMF supplies that
hypothesis, and it does not prove the reverse implication.

\section{Source question}

The source is F. Di Plinio and Y. Ou, \emph{Banach-valued multilinear singular
integrals}, arXiv:1506.05827v3, Remark 3.7 on printed page 14
\cite{DiPlinioOu}.  The authors ask whether dyadic RMF and nontangential RMF
are equivalent.

\begin{center}
\includegraphics[width=\textwidth]{figures/open_problem_crop.png}
\end{center}

Fix the triple of Banach spaces and the trilinear identification used in the
source.  For a family $S\subset X_3$, write $\Rbd(S)$ for the corresponding
Rademacher bound.  Let $\PhiClass$ be the normalized grand class of kernels
from (3.1) of the source, and put
\[
 \M_{\PhiClass,\alpha}f(x)
 =\sup_{\phi\in\PhiClass}\Rbd\{f*\phi_t(y):|y-x|<\alpha t,\ t>0\}.
\]
For the canonical dyadic lattice $\mathcal D$, put
\[
 \M^\triangle f(x)=
 \Rbd\{\langle f\rangle_Q:Q\in\mathcal D,\ x\in Q\}.
\]
Multiplying the defining seminorm of $\PhiClass$ by a fixed constant changes
only harmless constants below.

\section{Idea of the proof}

A sharp cube average can be reached from a smooth cube average by letting the
smoothing scale tend to zero.  Differentiating in that scale and integrating
by parts moves the derivative to the boundary of the cube.  At relative scale
$2^{-m}$, the boundary term has total mass $O(2^{-m})$, while all of its smooth
averages lie in a cone of aperture $O(2^m)$ above every point of the cube.
Thus
\[
 \M^\triangle f
 \lesssim \M_{\PhiClass,C}
       +\sum_{m\geq0}2^{-m}\M_{\PhiClass,C2^m}f.
\]
The displayed pointwise estimate is unconditional.  The conditional
dependency enters only when taking $L^p$ norms and summing the aperture
constants.

\section{The reduction}

For $1<p<\infty$ and $\alpha\geq1$, denote the aperture norm by
\[
 A_p(\alpha)=
 \|\M_{\PhiClass,\alpha}\|_{L^p(\mathbb R^d;X_3)\to L^p(\mathbb R^d)}.
\]

\begin{theorem}[summable-aperture criterion]\label{thm:main}
There are dimensional constants $c_d,C_d>0$ such that, for every Schwartz
function $f:\mathbb R^d\to X_3$,
\begin{equation}\label{eq:pointwise}
 \M^\triangle f(x)
 \leq C_d\left(
 \M_{\PhiClass,c_d}f(x)
 +\sum_{m=0}^{\infty}2^{-m}
       \M_{\PhiClass,c_d2^m}f(x)\right)
\end{equation}
for almost every $x$.  Consequently, if
\begin{equation}\label{eq:summability}
 A_p(c_d)+\sum_{m=0}^{\infty}2^{-m}A_p(c_d2^m)<\infty,
\end{equation}
then $X_3$ has dyadic $\mathrm{RMF}_p$.  In particular, it is enough that
\[
 A_p(\alpha)\leq C\alpha^\theta\qquad(\alpha\geq1)
\]
for some $\theta<1$.
\end{theorem}

\begin{proof}
Choose $\eta\in C_c^\infty(\mathbb R^d)$ with $\int\eta=1$.  After one
fixed normalization, $\eta$ and
\[
 \phi_i(u)=u_i\eta(u),\qquad i=1,\ldots,d,
\]
belong to the grand kernel class.  If $Q$ is a cube of side length $\ell$
and $r>0$, define
\[
 H_Q(r)=\frac1{|Q|}\int_Q(f*\eta_r)(y)\,dy.
\]
For Schwartz $f$, $H_Q(r)\to\langle f\rangle_Q$ in $X_3$ as $r\downarrow0$.
Moreover,
\[
 \partial_r\eta_r=-\sum_{i=1}^d\partial_i(\phi_i)_r.
\]
The divergence theorem, valid here for the smooth $X_3$-valued convolutions,
therefore gives
\begin{equation}\label{eq:derivative}
 H'_Q(r)=-\frac1{|Q|}\sum_{i=1}^d
 \int_{\partial Q}(f*(\phi_i)_r)(y)n_i(y)\,d\sigma(y).
\end{equation}
It follows that
\begin{equation}\label{eq:decomposition}
 \langle f\rangle_Q=H_Q(\ell)+\sum_{m=0}^{\infty}B_{m,Q},
 \qquad
 B_{m,Q}=-\int_{2^{-m-1}\ell}^{2^{-m}\ell}H'_Q(r)\,dr.
\end{equation}

We use the standard facts that Rademacher bounds are unchanged by passing to
the closed absolutely convex hull, are preserved by normalized Bochner
averaging, and are subadditive under Minkowski sums; these are also used in
the source (notably in its proof of Lemma 3.2).

Fix $x$ and a finite collection $Q_1,\ldots,Q_N\in\mathcal D$ containing
$x$.  Since $|y-x|\leq\sqrt d\,\ell(Q_j)$ for $y\in Q_j$, normalized
averaging in the definition of $H_{Q_j}(\ell(Q_j))$ yields
\begin{equation}\label{eq:mainterm}
 \Rbd\{H_{Q_j}(\ell(Q_j)):1\leq j\leq N\}
 \leq C_d\M_{\PhiClass,c_d}f(x).
\end{equation}

For $r\in(2^{-m-1}\ell(Q_j),2^{-m}\ell(Q_j)]$ and
$y\in\partial Q_j$, one has
\[
 |y-x|\leq\sqrt d\,\ell(Q_j)\leq 2^{m+1}\sqrt d\,r.
\]
On the other hand, the total variation of the normalized boundary--scale
measure in \eqref{eq:derivative}, restricted to this interval of $r$, is at
most
\[
 \frac{|\partial Q_j|}{|Q_j|},2^{-m-1}\ell(Q_j)
 =d,2^{-m}.
\]
Splitting the $d$ kernel components and using preservation under Bochner
averaging therefore gives
\begin{equation}\label{eq:boundary}
 \Rbd\{B_{m,Q_j}:1\leq j\leq N\}
 \leq C_d2^{-m}\M_{\PhiClass,c_d2^m}f(x).
\end{equation}
Combining \eqref{eq:decomposition}--\eqref{eq:boundary}, using subadditivity,
and then letting the finite collection exhaust the dyadic cubes through $x$
proves \eqref{eq:pointwise}.  Minkowski's inequality and
\eqref{eq:summability} give
\[
 \|\M^\triangle f\|_p
 \leq C_d\left(A_p(c_d)+\sum_{m\geq0}2^{-m}A_p(c_d2^m)\right)
 \|f\|_{L^p(X_3)}.
\]
Density completes the $L^p$ statement.  If $A_p(\alpha)\leq C\alpha^\theta$
with $\theta<1$, the remaining series is geometric.
\end{proof}

\section{Conditional dependencies and limitations}

The boundary reconstruction and the pointwise estimate
\eqref{eq:pointwise} are complete.  The unproved dependency is precisely
\eqref{eq:summability}: ordinary nontangential RMF gives finiteness of each
fixed-aperture norm but no summable quantitative growth estimate is currently
derived here.  The condition cannot simply be replaced by aperture-uniform
boundedness without justification.  Already in the scalar case, for a
nonnegative compactly supported approximate identity and a test function
supported on a unit ball, the aperture-$\alpha$ maximal function stays bounded
below on a region of volume comparable to $\alpha^d$; hence its $L^p$ norm
has growth at least of order $\alpha^{d/p}$.

Even after proving \eqref{eq:summability}, this packet establishes only the
nontangential-to-dyadic implication.  The dyadic-to-nontangential implication
remains untouched.

The most important reviewer task is to check the passage from the boundary
formula \eqref{eq:derivative} to the simultaneous Rademacher-bound estimate
\eqref{eq:boundary}, including the normalization of the finitely many kernels
$u_i\eta(u)$ into the source's grand class.

\section{Verification and novelty status}

The proof was checked symbolically: differentiating
$r^{-d}\eta(z/r)$ gives
$-r^{-d-1}(d\eta+(z/r)\cdot\nabla\eta)$, exactly
$-\sum_i\partial_i(\phi_i)_r$.  The boundary mass on the $m$th dyadic scale
is $d2^{-m}$, and the required aperture is at most
$2^{m+1}\sqrt d$.  No numerical or finite-dimensional experiment is used.

A bounded search on 21 July 2026 covered the exact phrases
``nontangential Rademacher maximal function'', ``nontangential RMF'', and
``dyadic RMF and nontangential RMF'', together with arXiv:1506.05827 and the
later arXiv papers 1905.02139 and 1908.07233.  No later arXiv paper explicitly
settling the equivalence or stating this boundary/aperture reduction was
found.  This is not a proof of novelty.

\begin{thebibliography}{9}
\bibitem{DiPlinioOu}
F. Di Plinio and Y. Ou,
\emph{Banach-valued multilinear singular integrals},
Indiana Univ. Math. J. 67 (2018), no. 5, 1711--1763;
arXiv:1506.05827.
\end{thebibliography}

\end{document}
